% ============================================================
% P4-IDS Academic Poster — ENHANCED EXPLANATORY VERSION
% ============================================================
\documentclass[landscape]{article}
\usepackage[paperwidth=48in,paperheight=36in,margin=0.65in]{geometry}
\usepackage{graphicx}
\usepackage[table]{xcolor}
\usepackage{array}
\usepackage{booktabs}
\usepackage{url}
\usepackage{ragged2e}
\usepackage{tikz}
\usepackage[numbers]{natbib}
\usetikzlibrary{shapes.geometric, arrows.meta, positioning, calc}

\pagestyle{empty}
\setlength{\parindent}{0pt}
\setlength{\parskip}{0.16in}
\setlength{\fboxsep}{0.14in}
\setlength{\fboxrule}{1.5pt}

\definecolor{slublue}{RGB}{0,56,101}
\definecolor{slugold}{RGB}{198,146,20}
\definecolor{lightblue}{RGB}{224,237,250}
\definecolor{darkgray}{RGB}{35,35,35}
\definecolor{medgray}{RGB}{95,95,95}
\definecolor{okgreen}{RGB}{22,120,40}
\definecolor{badred}{RGB}{180,30,30}
\definecolor{blockbg}{RGB}{246,249,252}
\definecolor{softgold}{RGB}{255,248,224}

\newcommand{\BodyText}{\fontsize{25}{31}\selectfont}
\newcommand{\SmallText}{\fontsize{18}{22}\selectfont}

\newcommand{\block}[2]{%
\vspace{0.16in}
{\color{slublue}\fbox{\begin{minipage}{\dimexpr\linewidth-2\fboxsep-2\fboxrule\relax}
\colorbox{slublue}{\parbox{\dimexpr\linewidth-2\fboxsep\relax}{\color{white}\fontsize{32}{38}\selectfont\bfseries #1}}\\[-0.02in]
\colorbox{blockbg}{\parbox{\dimexpr\linewidth-2\fboxsep\relax}{\vspace{0.08in}\color{darkgray}\BodyText\justifying #2\vspace{0.08in}}}
\end{minipage}}}}

\newcommand{\alertblock}[2]{%
\vspace{0.16in}
{\color{slugold}\fbox{\begin{minipage}{\dimexpr\linewidth-2\fboxsep-2\fboxrule\relax}
\colorbox{slugold}{\parbox{\dimexpr\linewidth-2\fboxsep\relax}{\color{white}\fontsize{32}{38}\selectfont\bfseries #1}}\\[-0.02in]
\colorbox{softgold}{\parbox{\dimexpr\linewidth-2\fboxsep\relax}{\vspace{0.08in}\color{darkgray}\BodyText\justifying #2\vspace{0.08in}}}
\end{minipage}}}}

\newcommand{\metricbox}[3]{%
\fcolorbox{#3}{#3!12}{\begin{minipage}[c][1.55in][c]{3.25in}\centering
{\color{#3}\fontsize{44}{50}\selectfont\bfseries #1}\\[0.03in]
{\color{darkgray}\fontsize{18}{22}\selectfont #2}
\end{minipage}}}

\begin{document}

% ================= Header =================
\noindent
\colorbox{slublue}{%
\begin{minipage}[c][4.6in][c]{\dimexpr\linewidth-2\fboxsep\relax}
\begin{minipage}[c][3.4in][c]{0.14\linewidth}\centering
\includegraphics[width=3.4in,height=3.0in,keepaspectratio]{sluu.png}
\end{minipage}
\hfill
\begin{minipage}[c]{0.72\linewidth}\centering
{\color{white}\fontsize{120}{150}\selectfont\bfseries P4-Based Intrusion Detection System}\\[0.05in]
{\color{white}\fontsize{75}{85}\selectfont\bfseries with Rule Distillation for Data-Plane Security}\\[0.15in]
{\color{slugold}\fontsize{50}{56}\selectfont\bfseries Venugopal Rao Ponnamaneni \quad \& \quad Moditha Lekkala}\\[0.20in]
{\color{white}\fontsize{38}{44}\selectfont Department of Computer Science, Saint Louis University}
\end{minipage}
\hfill
\begin{minipage}[c][3.4in][c]{0.12\linewidth}\centering
\includegraphics[width=3.0in,height=3.0in,keepaspectratio]{qr.png}\\[-0.03in]
{\color{white}\fontsize{22}{26}\selectfont GitHub Repo}
\end{minipage}
\end{minipage}}

\vspace{0.3in}

% ================= Start Columns =================
\noindent
\begin{minipage}[t]{0.31\linewidth}
\block{Abstract}{
\textbf{Context:} High-speed Software-Defined Networks (SDN) require instantaneous security responses to mitigate modern cyber threats. 
\textbf{Need:} Current Intrusion Detection Systems (IDS) typically mirror traffic to external controllers, creating a "latency gap" where attacks succeed before a mitigation command is ever issued. 
\textbf{Task:} We designed a line-rate IDS that resides entirely within the data plane using the P4 language and BMv2 switches. 
\textbf{Object:} This poster details the transformation of Machine Learning models into hardware-compatible logic via rule distillation. 
\textbf{Findings:} Our implementation achieves 100\% precision and zero false positives on the CICIDS2017 dataset. 
\textbf{Meanings:} The results prove that complex ML-driven security can be offloaded to programmable hardware without sacrificing accuracy. 
\textbf{Perspectives:} Future work involves optimizing TCAM memory usage for deployment on Tofino-based hardware targets.
}

\alertblock{The Motivation: Closing the Latency Gap}{
The core challenge in network security is the delay between \textbf{detection} and \textbf{mitigation}. 
\begin{itemize}
    \item \textbf{Traditional Method:} Data Plane $\rightarrow$ Control Plane $\rightarrow$ CPU Analysis $\rightarrow$ Rule Update. This takes milliseconds to seconds.
    \item \textbf{Our Proposed Goal:} Detection happens \textit{at the switch} in microseconds, allowing the first malicious packet to be dropped before it exits the ingress pipeline.
\end{itemize}
}

\alertblock{The Problem: The "Latency Gap"}{
Traditional Intrusion Detection Systems (IDS) rely on \textbf{off-path analysis}, where traffic is mirrored to a central controller or CPU. This creates a critical delay:
\begin{itemize}
    \item \textbf{Slow Response:} Malicious packets often reach their destination before the IDS can issue a DROP command.
    \item \textbf{Controller Bottleneck:} High-speed links (100Gbps+) overwhelm software-based analyzers.
    \item \textbf{Resource Waste:} Constant data transfer between the data plane and control plane consumes massive bandwidth.
\end{itemize}
}

\block{Technical Foundation}{
\textbf{P4 Language:} Defines packet parsing, metadata, and match-action tables. \\
\textbf{BMv2:} Behavioral software switch target used for the P4 pipeline. \\
\textbf{TCAM:} Hardware memory enabling parallel ternary matching (value, mask, wildcard).
}

\block{System Workflow}{
The architecture follows a "Train Offline, Execute Online" philosophy to bypass traditional CPU bottlenecks.
\renewcommand{\arraystretch}{1.28}
\begin{center}
\begin{tabular}{>{\centering\arraybackslash}p{13.5cm}}
\rowcolor{softgold}\textbf{Offline Training and Feature Extraction} \\
\SmallText (Uses historical PCAP data to identify key malicious flow features) \\
$\Downarrow$ \\
\rowcolor{softgold}\textbf{Distill Model Logic into \texttt{rules.json}} \\
\SmallText (Converts Random Forest/Trees into Ternary Match-Action rules) \\
$\Downarrow$ \\
\rowcolor{lightblue}\textbf{Incoming Packets Enter BMv2 Switch} \\
\SmallText (Real-time traffic is parsed at the ingress pipeline) \\
$\Downarrow$ \\
\rowcolor{lightblue}\textbf{TCAM-Style Ternary Table Lookup} \\
\SmallText (Hardware-level comparison with zero software interrupts) \\
$\Downarrow$ \\
\textbf{Decision: DROP / FORWARD}
\end{tabular}
\end{center}
}
\end{minipage}\hfill
%
\begin{minipage}[t]{0.36\linewidth}
\block{Architecture: Control vs. Data Plane}{
In this implementation, the separation of concerns ensures both intelligence and speed:

\textbf{Control Plane (The "Brain"):}
\begin{itemize}
    \item \textbf{Function:} Operates offline or asynchronously to the traffic flow.
    \item \textbf{In this project:} It handles the heavy lifting—training the Machine Learning model on the CICIDS2017 dataset and "distilling" those complex patterns into simple Ternary match-action rules.
    \item \textbf{Output:} It pushes the \texttt{rules.json} entries into the switch tables.
\end{itemize}

\textbf{Data Plane (The "Brawn"):}
\begin{itemize}
    \item \textbf{Function:} Operates at line rate (nanoseconds) on every single packet.
    \item \textbf{In this project:} It uses the P4 pipeline to parse headers, aggregate flow features in registers, and perform a $O(1)$ TCAM lookup against the distilled rules.
    \item \textbf{Outcome:} It makes the final autonomous decision to \textbf{FORWARD} or \textbf{DROP} without waiting for the Controller.
\end{itemize}
}

\block{P4 Data-Plane Implementation}{
The P4 code acts as the programmable forwarding logic handling packets at line rate:
\begin{itemize}
    \item \textbf{Parser:} Extracts IP/TCP/UDP header fields.
    \item \textbf{Feature Aggregation:} Uses stateful registers to calculate packet sizes and flow durations over a sliding window.
    \item \textbf{Ingress Control:} Executes match-action logic with distilled ML rules to classify traffic in a single clock cycle.
    \item \textbf{Deparser:} Reconstructs headers for benign traffic to ensure seamless forwarding.
\end{itemize}
}

\block{ML Model: Centroid-Based Classification}{
To maintain line-rate performance, this project uses operations that can be executed efficiently in a switch pipeline without floating-point math.

\textbf{Bitwise XNOR:} Compares encoded input features against normal and attack centroids to measure binary similarity.

\textbf{Popcount:} Counts the total number of matching bits after the XNOR comparison to generate a "proximity score."

\textbf{Decision:} If the attack centroid has a higher similarity score than the normal centroid, the packet is flagged as malicious.

\textbf{Speed ($O(1)$ Complexity):} The switch compares input traffic against fixed centroids, so inference cost remains constant regardless of traffic volume.
}

\block{ML Model: Threshold point}{
\begin{minipage}{0.55\linewidth}
    \BodyText
    Traffic classes are represented as prototype bit-strings called \textbf{centroids}. The switch compares incoming features using:
    \begin{itemize}
        \item \textbf{Bitwise XNOR:} Measures bit-level alignment.
        \item \textbf{Popcount:} Calculates similarity scores.
    \end{itemize}
    This enables autonomous, line-rate mitigation by eliminating external controller bottlenecks.
\end{minipage}
\hfill
\begin{minipage}{0.4\linewidth}
    \centering
    \includegraphics[width=\linewidth,keepaspectratio]{random.png}\\
    {\SmallText\color{medgray}\textbf{Fig. 1:} ROC Curve (AUC = 0.900)\par}
\end{minipage}

\vspace{0.1in}
\textbf{Inference:} \textit{The high AUC of 0.90 indicates that our centroid-based logic effectively separates attack traffic from benign traffic with high sensitivity.}
}
\end{minipage}\hfill
%
\begin{minipage}[t]{0.31\linewidth}
\block{Knowledge Distillation \& TCAM Deployment}{
Decision-tree behavior is flattened into TCAM value-mask-action rules, reducing lookup complexity from $O(\text{depth})$ to $O(1)$.

\begin{center}
\resizebox{0.98\linewidth}{!}{%
\begin{tikzpicture}[
    font=\sffamily,
    tree_node/.style={rectangle, rounded corners, draw=slublue, fill=slublue!10, minimum width=4.5cm, minimum height=1.5cm, font=\bfseries\Large, align=center},
    leaf_drop/.style={rectangle, rounded corners, draw=badred, fill=badred!20, minimum width=3.2cm, minimum height=1.3cm, font=\bfseries\Large, align=center},
    leaf_fwd/.style={rectangle, rounded corners, draw=okgreen, fill=okgreen!20, minimum width=3.2cm, minimum height=1.3cm, font=\bfseries\Large, align=center},
    table_node/.style={rectangle, draw=darkgray, fill=white, minimum width=11.5cm, minimum height=1.6cm, font=\ttfamily\LARGE, align=center},
    arrow/.style={-Stealth, line width=6pt, slublue}
]
    \node[tree_node] (root) at (0,0) {pkt\_len $\le$ 100?};
    \node[tree_node] (L1) at (-3.5,-2.8) {psh = 0?};
    \node[tree_node] (R1) at (3.5,-2.8) {psh $>$ 0?};
    \node[leaf_drop] (D1) at (-5.0,-5.1) {DROP};
    \node[leaf_fwd] (F1) at (-2.0,-5.1) {FWD};
    \node[leaf_fwd] (F2) at (2.0,-5.1) {FWD};
    \node[leaf_drop] (D2) at (5.0,-5.1) {DROP};

    \draw[line width=3pt] (root) -- (L1); \draw[line width=3pt] (root) -- (R1);
    \draw[line width=3pt] (L1) -- (D1); \draw[line width=3pt] (L1) -- (F1);
    \draw[line width=3pt] (R1) -- (F2); \draw[line width=3pt] (R1) -- (D2);

    \draw[arrow] (6.4,-2.8) -- (11.4,-2.8);
    \node[font=\bfseries\huge, black] at (8.9,-1.85) {Flatten + Merge};
    \node[font=\bfseries\huge, black] at (8.9,-3.85) {$O(\text{depth}) \rightarrow O(1)$};

    \node[table_node, fill=slublue, text=white] (h) at (18.2,0.15) {\bfseries Value \qquad Mask \qquad Action};
    \node[table_node, below=0pt of h] (r1) {0x003C \quad 0xFFFF \quad \color{badred}\bfseries DROP};
    \node[table_node, below=0pt of r1] (r2) {0x01F4 \quad 0xFF00 \quad \color{okgreen}\bfseries FWD};
    \node[table_node, below=0pt of r2] (r3) {0x01E0 \quad 0xFFF0 \quad \color{okgreen}\bfseries FWD};
    \node[table_node, below=0pt of r3] (r4) {0x0000 \quad 0xFFFF \quad \color{badred}\bfseries DROP};
\end{tikzpicture}
}
\end{center}
\textbf{Learning:} \textit{Flattening the tree into TCAM rules allows us to bypass sequential "if-else" logic, making security decisions in a single hardware lookup cycle.}
}

\block{Mitigation Timeline}{
\centering
\includegraphics[width=0.98\linewidth,height=5.5in,keepaspectratio]{timegap.png}
{\SmallText\color{medgray}\textbf{Fig. 2:} Two-phase pipeline. Delay $\approx$ 4.8s.\par}
\vspace{0.05in}
\textbf{Learning:} \textit{The timeline highlights the latency window where attack packets are actively processed before the switch successfully transitions to a "Mitigation" state.}
}

\alertblock{Experimental Results}{
\textbf{Confusion Matrix --- 100 Test Flows}
\begin{center}
\renewcommand{\arraystretch}{1.45}
\begin{tabular}{c|c|c}
& \color{badred}\textbf{Predicted DROP} & \color{okgreen}\textbf{Predicted FWD} \\
\hline
\color{badred}\textbf{Actual DROP} & \cellcolor{okgreen!15}\textbf{36 (TP)} & \cellcolor{badred!12}\textbf{14 (FN)} \\
\hline
\color{okgreen}\textbf{Actual FWD} & \cellcolor{badred!12}\textbf{0 (FP)} & \cellcolor{okgreen!15}\textbf{50 (TN)} \\
\end{tabular}
\end{center}
\metricbox{86\%}{Accuracy}{slublue}\hspace{0.05in}
\metricbox{100\%}{Precision}{okgreen}\hspace{0.05in}
\metricbox{72\%}{Recall}{slugold}

\vspace{0.1in}
\textbf{Inference:} \textit{The 100\% Precision confirms that our system never drops legitimate traffic (Zero False Positives), ensuring network reliability while maintaining robust security.}
}


\block{Conclusion}{%
    This work demonstrates that \textbf{ML-grade network security can reside entirely inside a programmable switch}, eliminating the round-trip bottleneck of traditional IDS architectures.

    \vspace{0.08in}
    Our distilled TCAM rules achieve \textbf{100\% precision} and \textbf{zero false positives}, meaning no legitimate user traffic was ever disrupted — a critical requirement for production deployments.

    \vspace{0.08in}
    The $O(1)$ lookup replaces $O(\text{depth})$ tree traversal, enabling \textbf{autonomous, line-rate mitigation} at the ingress pipeline with no software overhead.
  }

\block{References}{
{\fontsize{18}{22}\selectfont
\renewcommand{\section}[2]{} % Remove the default "References" title
\nocite{*} % Cite everything in the .bib file
\bibliographystyle{plainnat}
\bibliography{capstone} % Reference capstone.bib
}}

\end{minipage}

\vspace*{\fill}
\noindent
\colorbox{slublue}{\parbox{\dimexpr\linewidth-2\fboxsep\relax}{\color{white}\fontsize{24}{28}\selectfont AI Capstone \enspace$\bullet$\enspace Saint Louis University \enspace$\bullet$\enspace P4 / BMv2 IDS \hfill \color{slugold}\textbf{Contact: modithalekkala@gmail.com}}}
\end{document}